% ══════════════════════════════════════════════════════════════
%                        CHAPTER 1
% ══════════════════════════════════════════════════════════════

\chapter{Project Context}

\section{Introduction}

Internships represent a fundamental component of higher education, serving as a bridge between academic learning and professional experience. They allow students to apply theoretical knowledge in real-world environments, develop practical skills, and build professional networks before entering the job market.

However, the process of finding, applying for, and managing internships remains largely fragmented and manual in many regions, particularly in Algeria. This chapter provides a comprehensive overview of the project context, examines the current state of internship management, analyzes existing solutions in the market, and identifies the limitations that motivate the development of Stag.io.

\section{General Context}

\subsection{The Internship Process in Higher Education}

In the higher education system, internships serve multiple critical purposes. For students, they provide hands-on experience in their field of study, help them discover career paths, and strengthen their employability upon graduation. For companies, internships represent an opportunity to identify and evaluate potential future employees while benefiting from fresh perspectives and up-to-date academic knowledge. For universities, internship programs are essential for maintaining the relevance of their curricula and ensuring their graduates are well-prepared for the professional world.

The typical internship lifecycle involves several stages:

\begin{enumerate}
    \item \textbf{Offer publication}: Companies identify their internship needs and publish offers through various channels.
    \item \textbf{Discovery and search}: Students search for offers that match their skills, academic background, and career interests.
    \item \textbf{Application}: Students submit their applications, typically including a CV and a cover letter.
    \item \textbf{Selection}: Companies review applications, conduct interviews, and select candidates.
    \item \textbf{Agreement validation}: The university validates the internship agreement (convention de stage), ensuring it meets academic requirements.
    \item \textbf{Internship execution}: The student completes the internship under the supervision of both the company and a university advisor.
    \item \textbf{Evaluation}: At the end of the internship, both the company and the university evaluate the student's performance.
\end{enumerate}

Each of these stages traditionally involves multiple stakeholders, substantial paperwork, and significant coordination effort.

\subsection{Digital Transformation of Internship Management}

The digital transformation has reshaped virtually every industry, and education is no exception. Over the past decade, numerous platforms have emerged to facilitate the connection between job seekers and employers. General-purpose job boards such as LinkedIn and Indeed have become major players in the recruitment landscape, offering vast databases of job and internship listings.

More recently, specialized platforms such as Handshake (in the United States) and JobTeaser (in Europe) have been developed specifically for the higher education ecosystem. These platforms integrate directly with university career centers and provide features tailored to early-career candidates and campus recruiting.

Key trends in the digital internship management space include:

\begin{itemize}
    \item \textbf{Intelligent matching}: Using algorithms to recommend offers based on a candidate's profile, moving beyond simple keyword search.
    \item \textbf{Real-time tracking}: Enabling all stakeholders to monitor application status and internship progress in real time.
    \item \textbf{Administrative automation}: Digitizing paperwork such as internship agreements, evaluations, and progress reports.
    \item \textbf{Multi-stakeholder platforms}: Creating unified environments where students, companies, and universities can collaborate on a single platform.
\end{itemize}

\subsection{The Algerian Context}

In Algeria, the internship management landscape presents unique challenges. The Ministry of Higher Education and Scientific Research has initiated a digital roadmap (sch\'ema directeur num\'erique) to modernize university services. Initiatives such as the ALUMNI platform aim to connect graduates with professional opportunities, including internships.

However, the current ecosystem remains fragmented:

\begin{itemize}
    \item \textbf{University portals} vary significantly in their capabilities, with some institutions offering basic administrative interfaces for managing internship paperwork while others have no digital tools at all.
    \item \textbf{No unified national platform} exists that centralizes internship offers across all sectors and enables structured collaboration between students, companies, and universities.
    \item \textbf{Existing job boards} such as Stapra Stage and EmploiPartner focus primarily on job listings, with limited features for the full internship lifecycle (agreement validation, real-time tracking, intelligent matching).
    \item \textbf{Students rely heavily on direct contact} (networking, LinkedIn, in-person visits to companies) because no single platform provides a reliable, centralized source of internship offers.
\end{itemize}

This gap in the market creates an opportunity for a purpose-built platform that addresses the specific needs of the Algerian higher education ecosystem while incorporating modern features found in international solutions.

\section{Problem Statement}

Based on our analysis of the current internship management landscape, we identify the following key problems:

\begin{enumerate}
    \item \textbf{Fragmentation of information}: Internship offers are scattered across multiple channels (company websites, social media, university bulletin boards, word of mouth), making it difficult for students to discover all relevant opportunities. There is no single source of truth for available internships.

    \item \textbf{Inefficient matching}: Existing platforms present offers in a generic chronological or alphabetical order. Students must manually evaluate each offer's relevance to their profile, which is time-consuming and leads to irrelevant applications that waste time for both students and companies.

    \item \textbf{Manual administrative processes}: The generation and validation of internship agreements (conventions de stage) still relies on paper-based workflows in most Algerian universities. This leads to delays, lost documents, and a lack of transparency in the approval process.

    \item \textbf{Lack of real-time visibility}: Once an application is submitted, students often have no way to track its status. Similarly, companies have no centralized tool to manage and compare applicants across multiple offers. Universities cannot easily monitor which students have secured internships.

    \item \textbf{Absence of feedback mechanisms}: After completing an internship, students have no structured way to share their experiences or rate companies. This lack of feedback prevents future students from making informed decisions and prevents companies from improving their internship programs.

    \item \textbf{No university integration}: Existing job boards and platforms operate independently from university systems. There is no mechanism to verify student enrollment, link students to their university, or enable universities to participate in the internship management process.
\end{enumerate}

These problems collectively result in a sub-optimal experience for all stakeholders: students struggle to find relevant internships, companies receive poorly matched applications, and universities lack the tools to effectively support and monitor their students' professional development.

\section{Presentation of Existing Solutions}

In this section, we examine the most prominent platforms currently used for internship search and management, analyzing their features, strengths, and limitations.

\subsection{LinkedIn}

\textbf{LinkedIn} is the world's largest professional networking platform, with over 1 billion members globally. It offers a comprehensive job and internship search functionality alongside its core networking features.

\textbf{Key features:}
\begin{itemize}
    \item Extensive job and internship listing database
    \item Professional profile creation and networking
    \item Company pages with reviews and insights
    \item InMail messaging for direct outreach to recruiters
    \item Skill endorsements and recommendations
    \item AI-powered job recommendations
\end{itemize}

\textbf{Limitations for internship management:}
\begin{itemize}
    \item Not specifically designed for the internship lifecycle
    \item No university integration or agreement validation
    \item No role-based access for different stakeholders
    \item Intelligent matching is basic and primarily keyword-based
    \item No administrative workflow for internship agreements
    \item Premium features require a paid subscription
\end{itemize}

\subsection{Handshake}

\textbf{Handshake} is a career platform specifically designed for the higher education ecosystem, widely used in the United States. It connects students directly with employers through university career centers.

\textbf{Key features:}
\begin{itemize}
    \item Direct integration with university career centers
    \item Employer vetting through university partnerships
    \item Virtual career fairs and events
    \item On-campus recruiting tools
    \item Targeted messaging between employers and students
    \item University-specific analytics and reporting
\end{itemize}

\textbf{Limitations for internship management:}
\begin{itemize}
    \item Primarily available in the United States and select countries
    \item Not available in Algeria or the North African region
    \item Limited support for multilingual environments
    \item No built-in internship agreement generation or validation workflow
    \item No intelligent scoring engine for offer-student matching
\end{itemize}

\subsection{Indeed}

\textbf{Indeed} is one of the world's largest job search engines, aggregating listings from thousands of company websites, job boards, and staffing agencies.

\textbf{Key features:}
\begin{itemize}
    \item Massive database of job and internship listings
    \item Simple, user-friendly search interface
    \item Resume upload and application tracking
    \item Company reviews and salary information
    \item Mobile application for job search on the go
\end{itemize}

\textbf{Limitations for internship management:}
\begin{itemize}
    \item Purely a job search engine with no internship-specific features
    \item No university integration whatsoever
    \item No role-based access control
    \item No real-time notifications or application status tracking from the employer side
    \item No administrative workflow for academic requirements
    \item Generic matching with no profile-based intelligent recommendations
\end{itemize}

\subsection{Stapra Stage}

\textbf{Stapra Stage} is a platform dedicated to internship and job offers in Algeria, representing one of the few localized solutions available in the Algerian market.

\textbf{Key features:}
\begin{itemize}
    \item Focused on the Algerian market
    \item Internship and job offer listings
    \item Basic search and filtering capabilities
    \item Company profiles
\end{itemize}

\textbf{Limitations for internship management:}
\begin{itemize}
    \item Limited feature set compared to international platforms
    \item No university integration or agreement validation
    \item No intelligent matching or recommendation engine
    \item No real-time notification system
    \item No multi-role architecture (student, company, university, admin)
    \item No review or feedback system
\end{itemize}

\subsection{University Portals}

Several Algerian universities have developed their own internal portals for managing internship-related administrative tasks.

\textbf{Key features:}
\begin{itemize}
    \item Administrative management of internship paperwork
    \item Student enrollment verification
    \item Internal communication between students and academic advisors
\end{itemize}

\textbf{Limitations:}
\begin{itemize}
    \item Each university operates its own isolated system with no interoperability
    \item No connection to the company ecosystem
    \item Limited to administrative tasks with no offer discovery features
    \item No intelligent matching or modern UX
    \item Often outdated interfaces with poor user experience
\end{itemize}

\section{Comparative Study}

The following table provides a comparative analysis of the existing solutions against the key features required for an effective internship management platform:

\begin{table}[H]
\centering
\caption{Comparative Study of Existing Solutions}
\label{tab:comparison}
\footnotesize
\begin{tabular}{|p{3.5cm}|c|c|c|c|c|c|}
\hline
\textbf{Feature} & \textbf{LinkedIn} & \textbf{Handshake} & \textbf{Indeed} & \textbf{Stapra} & \textbf{Uni Portals} & \textbf{Stag.io} \\
\hline
Offer listing & \checkmark & \checkmark & \checkmark & \checkmark & $\times$ & \checkmark \\
\hline
Intelligent matching & $\times$ & $\times$ & $\times$ & $\times$ & $\times$ & \checkmark \\
\hline
University integration & $\times$ & \checkmark & $\times$ & $\times$ & \checkmark & \checkmark \\
\hline
Agreement validation & $\times$ & $\times$ & $\times$ & $\times$ & Partial & \checkmark \\
\hline
Multi-role access & $\times$ & Partial & $\times$ & $\times$ & $\times$ & \checkmark \\
\hline
Real-time notifications & $\times$ & \checkmark & $\times$ & $\times$ & $\times$ & \checkmark \\
\hline
Review system & \checkmark & $\times$ & \checkmark & $\times$ & $\times$ & \checkmark \\
\hline
CV generation & $\times$ & $\times$ & $\times$ & $\times$ & $\times$ & \checkmark \\
\hline
AI chatbot & $\times$ & $\times$ & $\times$ & $\times$ & $\times$ & \checkmark \\
\hline
Bilingual (EN/FR) & \checkmark & $\times$ & \checkmark & $\times$ & $\times$ & \checkmark \\
\hline
Available in Algeria & \checkmark & $\times$ & \checkmark & \checkmark & \checkmark & \checkmark \\
\hline
Free to use & Partial & \checkmark & \checkmark & \checkmark & \checkmark & \checkmark \\
\hline
\end{tabular}
\end{table}

This comparison clearly shows that no existing solution provides a complete, integrated platform that addresses all the needs of the internship management lifecycle. Each existing solution covers only a subset of the required features, and none offers intelligent matching combined with university integration and administrative workflow automation.

\section{Proposed Solution}

To fill the gap identified in our comparative study, we propose \textbf{Stag.io} --- a modern, full-stack web platform that unifies the internship management experience for all stakeholders. Stag.io distinguishes itself from existing solutions through several key innovations:

\begin{itemize}
    \item \textbf{SmartMatch Engine}: A proprietary, zero-dependency scoring algorithm that ranks offers by relevance to each student's profile across four dimensions (skills, department, location, title), producing real-time results in under 100\,ms.
    \item \textbf{Multi-stakeholder architecture}: A role-based system with five distinct roles (student, company, university, administrator, super administrator), each with tailored dashboards and capabilities.
    \item \textbf{End-to-end lifecycle coverage}: From offer publication and intelligent discovery to application, agreement validation, and post-internship reviews --- all within a single platform.
    \item \textbf{University-linked authentication}: Students register using their university email domain, automatically linking them to their institution and enabling university-level oversight.
    \item \textbf{Modern technology stack}: Built with Next.js 16, React 19, Express 5, PostgreSQL, and Socket.IO, ensuring performance, scalability, and a premium user experience.
\end{itemize}

The detailed presentation of the project, including actor identification, functional and non-functional requirements, and use case modeling, will be covered in the following chapter.

\section{Conclusion}

In this chapter, we have established the context for the Stag.io project by examining the internship management landscape both globally and within Algeria. We identified the key problems that students, companies, and universities face in the current process: fragmentation of information, inefficient matching, manual administrative procedures, lack of real-time visibility, and absence of structured feedback mechanisms.

Our analysis of existing solutions --- including LinkedIn, Handshake, Indeed, Stapra Stage, and university portals --- revealed that no single platform provides a complete, integrated solution for all stakeholders. This gap motivates the development of Stag.io, a platform designed to address these limitations through intelligent matching, multi-stakeholder collaboration, and end-to-end lifecycle coverage.

In the next chapter, we will present the preliminary study and requirements specification, including the identification of actors, functional and non-functional requirements, and use case modeling.
